\documentclass[11pt,a4paper]{article}

\usepackage[utf8]{inputenc}

\usepackage[T1]{fontenc}

\usepackage[polish]{babel}

\usepackage{lmodern}

\usepackage{geometry}

\geometry{margin=27mm}

\usepackage[hidelinks]{hyperref}

\title{Podłużny transfer uczenia regulacji poza grę}

\author{Mateusz Pomianek}

\date{Szkic naukowy, wersja 0.1.3}

\begin{document}

\maketitle

\begin{abstract}

Praca przedstawia odroczoną koncepcję badania transferu umiejętności regulacji nabywanych podczas gry do sytuacji bez gry. Proponuje rozróżnienie poprawy wykonania zadania treningowego od uogólnienia umiejętności na nowe zadanie. Pierwszorzędowy punkt końcowy i estymand wymagają określenia w przyszłym protokole. Nie zgromadzono materiału empirycznego ani nie przedstawiono wyników.

\end{abstract}

\section{Wprowadzenie}

Kontekst obserwacji i wieloźródłowość ocen są kluczowe w interpretacji zachowania i zjawisk neurorozwojowych. Zagadnienie Podłużny transfer uczenia regulacji poza grę wymaga odróżnienia badanego konstruktu od jego operacjonalizacji i kontekstowych źródeł zmienności. Trzeba porównać możliwe wyjaśnienia i ustalić granice interpretacji dowodów.

\section{Cel i pytania badawcze}

Planowane pytanie badawcze dotyczy tego, czy umiejętności regulacji ćwiczone w grze utrzymują się w czasie i przenoszą na nowe zadania wykonywane bez gry. Hipotezy kierunkowe, okres obserwacji oraz pierwszorzędowy punkt końcowy wymagają ustalenia przed rozpoczęciem gromadzenia danych.

\section{Materiał i metody}

Pierwszorzędowy wskaźnik: Samodzielne wykonanie nowego zadania regulacyjnego bez gry, jako kandydat na wynik główny. Jest to propozycja, nie zatwierdzony punkt końcowy.


Wstępna koncepcja zakłada pomiar wyjściowy, okres ćwiczeń z grą oraz odroczoną ocenę transferu w zadaniu wykonywanym bez gry, z odpowiednio dobranym warunkiem porównawczym. Samodzielne wykonanie nowego zadania regulacyjnego stanowi jedynie kandydaturę na wynik główny, a nie zatwierdzony wskaźnik. Operacjonalizacja transferu musi odróżniać wykonanie wyuczonego zadania od generalizacji umiejętności. Projekt nie ma potwierdzonego protokołu ani danych. Przed badaniem z udziałem ludzi należy określić kryteria doboru próby, liczebność i terminy pomiarów, procedurę zgody, wymogi etyczne oraz ochronę danych.

\section{Plan analizy}

Rozważa się porównanie zmian między warunkami badania z uwzględnieniem poziomu wyjściowego oraz brakujących obserwacji. Wybór modelu, estymandu, sposobu postępowania z brakami danych, kryteriów jakości i analiz wrażliwości wymaga uzasadnienia w przyszłym protokole przed analizą danych. Nie oszacowano żadnych efektów.

Zamrożenie estymandu i reguł analizy musi nastąpić przed dostępem do danych.

\section{Wyniki}

W niniejszym pliku nie umieszczono wyników empirycznych, statystyk ani nieprzeprowadzonych analiz. Sekcja wymaga uzupełnienia po uzyskaniu i weryfikacji materiału dowodowego.

\section{Dyskusja i ograniczenia}

Ograniczenia obejmują możliwe błędy pomiaru, dobór próby, alternatywne mechanizmy i ograniczoną możliwość uogólniania. Korelacja nie dowodzi przyczynowości; poprawność techniczna nie jest dowodem znaczenia klinicznego.

\section{Wnioski}

Zdefiniowano ramy opracowania, ale nie sformułowano wniosków wykraczających poza dostępny stan dowodów.

\section*{Bibliografia}

Weryfikacja literatury i przypisanie cytowań do stwierdzeń są kolejnym etapem prac.

\end{document}

